\scanpage{68}
% archive foliation: 67
[Cela est le cas en particulier dans le cas de l'exemple I, \deleted{$S'$ est} \inserted{lorsque $S'$ est} fini et loc. libre sur $S$, et l'exemple II, \inserted{lorsque} $W$ \deleted{\unsure{est}} est affine sur $S$].

\uline{Proposition} Dans ces conditions, si \deleted{$\mathfrak X$ est} $\mathfrak X$ est un schéma \uline{affine} sur $\underline W(S)$, $W^a(\mathfrak X)$ existe et est un schéma \uline{affine} sur $S$.

\uline{Démonstration} \deleted{\unsure{Soit}\ $\mathfrak X$} \inserted{On peut supposer $S$ affine.} Si $\mathfrak X=\operatorname{Spec}(\unsure{C})$ où $\unsure{C}$ est une algèbre de type fini \deleted{de} sur $\underline W(A)$, engendrée par $r$ éléments, il résulte de suite du \unsure{\S\ N° 1} que $W^a(\mathfrak X)$ existe et est $\simeq$ un sous-préschéma fermé de $W^r$ [défini par des ``équations algébriques'' à coeffs. \inserted{de} dans $\Gamma(W/S)$]. Comme $W$ est affine, il en est de $=$ de \deleted{$\unsure{W^r}$} $W^r$ donc aussi de $W^a(\mathfrak X)$ (\unsure{cqf}). On obtient d'ailleurs

\uline{Cor.} \deleted{Si \unsure{S aff.}\ \unsure{l\textquotesingle anneau de}} \inserted{Supposons $S$ affine.} Si $\unsure{\ell|C}$ est engendré par $r$ générateurs sur celui de $W(S)$, et celui de $W$ par $\unsure{v}$ génér. sur celui de $S$, alors
